% SPDX-License-Identifier: LPPL-1.3c
% Copyright (C) 2026 Christophe Jorssen
% Author and Current Maintainer: Christophe Jorssen <christophe.jorssen@gmail.com>
% LPPL maintenance status: maintained
% This file is part of luafemm. See LICENSE and LICENSES.md for its terms.
\section{Transient fields and controlled motion}
\label{sec:time}

A harmonic field describes one frequency after transients have disappeared.
A time calculation instead starts with a specified field and advances it through
successive instants. An impressed current can change at each step; conducting
regions oppose the change through induced currents. The previous potential is
part of the calculation, not merely a convenient initial guess for Newton.

The supported time formulation uses real planar or axisymmetric potentials,
reversible linear or nonlinear magnetic materials, permanent magnets, resolved
conductors and current-fed series windings. Set |depth| for a planar time model.
Full-circle P1 bands additionally permit prescribed motion and partitioned
single-axis rotational dynamics. External electrical networks, solid winding
terminals, ideal circuits, axial H diffusion, native/Fourier moving gaps and
history-dependent hysteresis are outside this time formulation. Harmonic loss
angles, effective wires and conducting homogenised sheets cannot supply a
time-domain material law. Resolve the conducting geometry instead.

\subsection{A conductor when a current is switched on}
Start with a coarse spatial mesh and a moderate number of time steps. Establish
the geometry, signs and characteristic time scale first. Then refine both the
mesh and the time step; one cannot compensate for the other.

The circular conductor below is initially unmagnetised. The two stranded source
regions turn on at the first positive time. They retain their impressed current
density without resolved eddy currents. The central conductor uses
|current constraint=zero|: its axial electric field adjusts so that its total
section current is zero. Declare disconnected isolated conductors as separate
regions if each needs its own zero-current constraint.

\luafemmexamplebegin{transient-shield}
\begin{codeexample}[width=0pt]
\tikzset{femm/materials/shield metal/.style={mu r=1,conductivity=5.8e7}}
\begin{tikzpicture}[
  femm/problem={xmin=-9,xmax=9,ymin=-7,ymax=7,mesh size=1,depth=10},
  femm/state={name=shield,initial=zero},scale=2]
  \filldraw[fill=orange!15,
    femm/region={material=shield metal,current constraint=zero}]
    (0,0) circle[radius=3];
  \filldraw[fill=red!15,
    femm/region={material=copper,conductor=stranded,current density=1e7}]
    (-6,-3) rectangle (-5,3);
  \filldraw[fill=blue!15,
    femm/region={material=copper,conductor=stranded,current density=-1e7}]
    (5,-3) rectangle (6,3);
  \node[above] at (0,3) {Initially unmagnetised conductor};
\end{tikzpicture}
\end{codeexample}

The order is deliberate: declare materials and regions, then capture the initial
state at the end of the picture. |initial=zero| requires an unsolved model and
homogeneous prescribed-potential boundaries. It sets the initial potential to
zero; it does not assert static equilibrium with the declared sources. Use
|initial=equilibrium|, the default, to solve a static initial condition instead.
Do not request a field solve before capturing a zero initial condition.

Each new picture below takes one step from the named state. At its end,
|femm/state| replaces that name with the newly converged snapshot. Previously
derived models, plots and profiles keep their own data. The displayed times are
0.2 and 2.4 milliseconds; the intervening steps still perform their solves.
The red symbols show induced currents normal to the drawing plane.

\begin{codeexample}[width=0pt]
\def\TimeResults{}
\foreach \Step in {1,...,12}{%
  \begin{tikzpicture}[
    femm/problem={previous state=shield,time step=.0002},
    femm/state={name=shield},scale=2]
    \femmselect[name=shield-body,material=shield metal]
    \femmmeasure[name=shield-\Step,quantity=integrals,selection=shield-body]
    \ifnum\Step=1
      \pic {femm geometry}; \pic[femm/lines=13] {femm field};
      \pic[femm/currents={material=shield metal,spacing=.8,
        reference current=1e7}] {femm eddy currents};
      \node[above] at (0,7) {$t=0.2\,\mathrm{ms}$};
    \fi
    \ifnum\Step=12
      \pic {femm geometry}; \pic[femm/lines=13] {femm field};
      \node[above] at (0,7) {$t=2.4\,\mathrm{ms}$};
    \fi
  \end{tikzpicture}%
  \xdef\TimeResults{\TimeResults,shield-\Step}%
}
\end{codeexample}

The result names form an ordinary PGFPlots data source. No plot is a second
simulation: each loss value belongs to the field that was just solved.
The selection restricts this curve to the shield, excluding winding heating.

\begin{codeexample}[width=0pt]
\begin{tikzpicture}
\begin{axis}[width=10cm,height=5.5cm,grid=major,
  xlabel={Time (ms)},ylabel={Shield Joule loss (W)},ymin=0,
  x filter/.expression={1000*x}]
  \femmresultplot[results=\TimeResults,component=joule_loss,abscissa=time,
    plot options={teal!70!black,thick,mark=*}]
\end{axis}
\end{tikzpicture}
\end{codeexample}

\begin{codeexample}[width=0pt]
\femmstate[name=shield-saved,file=\jobname-shield.lft]
\femmrestorestate[name=shield-restart,file=\jobname-shield.lft]
\end{codeexample}
\luafemmexampleend{transient-shield}

\subsection{State ownership and the order of operations}
\label{sec:time-keys}
\begin{key}{/tikz/femm/state=\meta{options}}
Capture the originating picture's model at its end, using |/femm/state| keys.
It remains bound to that picture even if a nested picture changes the current
model. Set |name| explicitly. Normal PGF styles may supply shared capture options.
\end{key}
\begin{command}{\femmstate\opt{\oarg{options}}}
Capture the current or named model immediately. A static model is solved on
demand; a time model must converge before its new state is registered. Options
belong to |/femm/state|. Reusing a state name replaces its owned snapshot.
\end{command}
\begin{key}{/femm/state/name=\meta{name} (initially empty)}
Required document-wide state identifier. A derived step copies its previous
potential and declarations; replacing this name never changes that step.
\end{key}
\begin{key}{/femm/state/model=\meta{name} (initially empty)}
Optional previously named problem. Empty selects the current model. The picture
key supplies its own model automatically. This key applies to capture only.
\end{key}
\begin{key}{/femm/state/initial=\meta{choice} (initially equilibrium)}
|equilibrium| solves an unsolved static initial problem. |zero| creates a zero
potential in an unsolved problem with homogeneous Dirichlet data. Existing
solved time fields are captured without changing their history; |zero| cannot
reset them. A frozen/incremental response cannot become a time state or vice versa.
\end{key}
\begin{key}{/femm/state/time=\meta{seconds} (initially empty)}
Initial absolute time, finite and nonnegative; empty means zero for an initial
state and the computed time for a solved step. A capture cannot relabel a solved
step's clock. Native answer import has its own clock key below.
\end{key}
\begin{key}{/femm/state/file=\meta{filename} (initially empty)}
Optionally write the captured state as passive data, conventionally with an
|.lft| extension. Relative names use the TeX output directory. The record owns
geometry, mesh, material/source data, present and previous potentials, clock and
optional motion template. It contains no executable Lua or TeX. Replacement is
atomic, and an incompatible version or damaged checksum is an error.
An |.lfc| cache accelerates an already declared calculation; an |.lft| file
supplies the state needed to restart one. Keep the source document as well.
\end{key}
\begin{command}{\femmrestorestate\opt{\oarg{options}}}
Read a saved state with |name| and |file| from |/femm/state|. No solve or time step
is performed. Only these two keys apply. A failed read leaves the existing named
state intact. Derive a step with |previous state| or display it with |state|.
\end{command}
\begin{key}{/femm/problem/state=\meta{name} (initially empty)}
Display or measure a saved snapshot without advancing time. The snapshot owns
its physical and numerical settings. Use ordinary TikZ transformations and
drawing styles for presentation. Excitation and motion overrides are errors.
\end{key}
\begin{key}{/femm/problem/previous state=\meta{name} (initially empty)}
Construct a new time step from this state. It is mutually exclusive with |state|,
|operating point|, |sliding band| and file import. The inherited geometry, material
laws and spatial mesh cannot be redeclared. Existing numerical solver-policy
keys may override the linear solver for the new step.
\end{key}
\begin{key}{/femm/problem/time step=\meta{seconds} (initially 0)}
Required positive finite step size with |previous state|. The new clock is
$t_{n+1}=t_n+\Delta t$. All impressed currents and boundary data belong to this
new time. A nonzero value without |previous state| is an error. The method is
backward Euler, first order in time; refine the step to assess numerical damping.
\end{key}
\begin{key}{/femm/problem/source scale=\meta{number} (initially empty)}
Multiply the initial state's non-circuit region current densities by this real,
dimensionless factor. Empty retains the preceding factor (initially one).
It does not scale magnetisation or named circuit currents. For winding currents,
reuse |/femm/problem/circuit currents={name=value,...}|, in amperes. Unmentioned
circuits retain their previous currents. Phase keys do not apply to time states.
\end{key}
\begin{key}{/femm/problem/boundary scale=\meta{number} (initially empty)}
Scale initial affine boundary loads and potentials, retaining Robin coefficients
and geometry. Empty retains the preceding factor (initially one). A scale of
zero homogenises these data; it does not change the boundary type. Automatic
static gauges are unsupported in time diffusion: provide an explicit potential
reference. The ordinary homogeneous exterior Dirichlet default meets this rule.
The circular open multipole approximation is also excluded: its static audit
does not account for induced return currents. Use a finite exterior domain and
check sensitivity to its extent. Explicit Robin data remain mathematical
boundary conditions, without an inferred open-exterior interpretation.
\end{key}

Spatial profiles retain the instantaneous |B|, |H|, |Az| and current components
of their source step. |J| and |Jeddy| use A/m$^2$; time profiles require
|representation=snapshot| and |phase=0|. The existing eddy-current pic draws
instantaneous normal symbols or axisymmetric loops; its harmonic phase setting
has no effect on a time snapshot. Planar B lines remain contours of $A_z$;
axisymmetric lines remain contours of full flux.

Regional measurements add |time| and |time_step| in seconds and |joule_loss| in
watts. This loss is instantaneous $\int J^2/\sigma\,dV$, over the selected model
volume, with winding fill included; it is not a harmonic cycle average.
|current| includes induced current. Maxwell stress and conductor Lorentz
measurements use the instantaneous field/current. Circuit measurements add
|induced_voltage| and use $V=RI+(\lambda_{n+1}-\lambda_n)/\Delta t$ when resistance
and closed-winding linkage are available. A user-specified winding length may
include end turns outside the field volume; its terminal resistance loss then
differs from the model-volume integral. No transient secant inductance is reported.

\subsection{Prescribed rotation and mechanical response}
\label{sec:time-motion}
\begin{key}{/femm/state/sliding band=\meta{name} (initially empty)}
Attach the original full-circle P1 template to an initial positioned model.
Its provenance must match the model's band. This stores the reference geometry
and material-node identities needed for later motion. A position alone does not
contain enough information to infer an arbitrary new motion template.
\end{key}
\begin{key}{/femm/state/speed=\meta{number} (initially empty)}
Initial angular speed in radians per second; empty means zero. Requires a motion
template and applies to an initial state only. Positive speed is counterclockwise.
\end{key}
\begin{key}{/femm/problem/motion=\meta{choice} (initially fixed)}
|fixed| holds geometry and sets the retained mechanical speed to zero.
|prescribed| uses |rotor angle| (degrees) and |rotor offset x/y| (model units) at
the new time; their ordinary initial values are zero. Supply unwrapped angles
when the derived speed matters. |dynamic| advances one rotational mechanical
degree of freedom; offsets remain fixed and position overrides are errors.
Moving steps require the initial state's attached P1 template. Only the
nonconducting air band is retriangulated. Conducting bulk elements retain node
identities, so their previous potential follows the material.
\end{key}
\begin{key}{/femm/problem/inertia=\meta{number} (initially 0)}
Required positive moment of inertia in kg\,m$^2$ for |motion=dynamic|. There is no
guessed physical inertia. It refers to the rotation axis through the band's
possibly displaced rotor centre. Use a PGF problem style to share it over steps.
\end{key}
\begin{key}{/femm/problem/load torque=\meta{number} (initially 0)}
Signed external opposing torque in N\,m. A positive value subtracts from the
electromagnetic torque, irrespective of the current speed sign. This is not an
automatic dry-friction law. Requires |motion=dynamic|.
\end{key}
\begin{key}{/femm/problem/damping=\meta{number} (initially 0)}
Nonnegative viscous coefficient in N\,m\,s per radian; its torque is $-c\omega$.
Requires |motion=dynamic|. No other mechanical load is inferred.
\end{key}

Here a reluctance rotor accelerates under magnetic torque. Its magnetic material
is nonconducting in this example, so the motion-induced terminal linkage and
mechanical response can be inspected separately from eddy-current braking.
Adding a resolved rotor conductivity activates the material time derivative.

\luafemmexamplebegin{transient-motion}
\begin{codeexample}[width=0pt]
\tikzset{femm/materials/rotor iron/.style={mu r=200}}
\begin{tikzpicture}[
  femm/problem={xmin=-8,xmax=8,ymin=-8,ymax=8,mesh size=1,depth=10},
  femm/sliding band={name=time-band,inner radius=3,outer radius=4,
    segments=48,layers=1},scale=2]
  \filldraw[fill=gray!20,femm/region={material=rotor iron}]
    (0,0) ellipse[x radius=2.6,y radius=1.2];
  \draw[femm/region={material=copper,conductor=stranded,current density=5e7}]
    (-6,-3) rectangle (-5,3);
  \draw[femm/region={material=copper,conductor=stranded,current density=-5e7}]
    (5,-3) rectangle (6,3);
\end{tikzpicture}
\begin{tikzpicture}[
  femm/problem={sliding band=time-band,rotor angle=25},
  femm/state={name=rotor-time,sliding band=time-band,speed=0},scale=2]
  \pic {femm geometry}; \pic[femm/lines=13] {femm field};
  \node[above] at (0,8) {$t=0$};
\end{tikzpicture}
\end{codeexample}

\begin{codeexample}[width=0pt]
\def\MotionResults{}
\foreach \Step in {1,...,20}{%
  \begin{tikzpicture}[
    femm/problem={previous state=rotor-time,time step=.001,
      motion=dynamic,inertia=1e-8,damping=1e-7},
    femm/state={name=rotor-time},scale=2]
    \femmmeasure[name=motion-\Step,quantity=force,method=band]
    \ifnum\Step=20
      \pic {femm geometry}; \pic[femm/lines=13] {femm field};
      \node[above] at (0,8) {$t=20\,\mathrm{ms}$};
    \fi
  \end{tikzpicture}%
  \xdef\MotionResults{\MotionResults,motion-\Step}%
}
\end{codeexample}

The saved measurements expose |angle| in degrees and |speed| in radians per
second. A time-history plot uses the existing result-plot interface.

\begin{codeexample}[width=0pt]
\begin{tikzpicture}
\begin{axis}[width=8cm,height=5cm,grid=major,
  xlabel={Time (ms)},ylabel={Rotor angle (degrees)},
  x filter/.expression={1000*x}]
  \femmresultplot[results=\MotionResults,abscissa=time,component=angle,
    plot options={blue!65!black,thick,mark=*}]
\end{axis}
\end{tikzpicture}
\femmstate[name=rotor-saved,file=\jobname-rotor.lft]
\end{codeexample}
\luafemmexampleend{transient-motion}

At each step the previous field supplies the band torque $T_n$. The partitioned
mechanical update is
\[
 \omega_{n+1}=\frac{I\omega_n+\Delta t(T_n-T_{\rm load})}{I+c\Delta t},
 \qquad \theta_{n+1}=\theta_n+\Delta t\,\omega_{n+1},
\]
where the second equation uses radians. The magnetic field is then solved on
the new position. This is a staggered first-order coupling, not a monolithic
implicit dynamics solve. Halve the step and compare both angle and torque;
backward-Euler diffusion stability does not certify a chosen mechanical step.
For imposed motion, compare the same physical positions at the same times.

\subsection{Time snapshots and FEMM 5}
\label{sec:time-exchange}
\begin{key}{/femm/import/time=\meta{seconds} (initially empty)}
Required absolute time for a supported transient |format=ans| import. It must be
finite and at least the file's positive |[DeltaT]|. Native answers store the step
size and previous potential but no absolute clock. This key is invalid for an
ordinary static/harmonic answer. The previous time is |time-DeltaT|.
\end{key}

The native exchange subset covers fixed planar real time snapshots, with no
AGE blocks or isolated-conductor voltage constraints. Current-fed stranded
windings and unconstrained resolved conductors are supported. Export preserves
the actual P1 mesh and both nodal potentials. These are FEMM 5 time records;
FEMM 4 does not supply the corresponding transient interpretation.

The following compact interchange check removes an impressed source from a
conducting square. The conductor has zero imposed longitudinal electric field,
not a prescribed zero total current. The two field pictures use the exported
and reimported snapshots, respectively.

\luafemmexamplebegin{transient-exchange}
\begin{codeexample}[width=0pt]
\tikzset{femm/materials/exchange metal/.style={mu r=2,conductivity=5e6}}
\begin{tikzpicture}[
  femm/problem={xmin=-3,xmax=3,ymin=-3,ymax=3,mesh size=.6,depth=10},
  femm/state={name=exchange-initial}]
  \draw[femm/region={material=exchange metal,current density=1e7}]
    (-2,-2) rectangle (2,2);
\end{tikzpicture}
\begin{tikzpicture}[
  femm/problem={previous state=exchange-initial,time step=.001,source scale=0},
  scale=4]
  \pic {femm geometry}; \pic[femm/lines=9] {femm field};
  \femmexport[format=ans,file=\jobname-time.ans]
\end{tikzpicture}
\begin{tikzpicture}[
  femm/import={file=\jobname-time.ans,time=.001},scale=4]
  \pic {femm geometry}; \pic[femm/lines=9] {femm field};
  \femmstate[name=imported-time,file=\jobname-imported.lft]
\end{tikzpicture}
\end{codeexample}
\luafemmexampleend{transient-exchange}

Import reads passive answer data only. A transient answer owns |DeltaT| and its
previous potential; a sibling solver-control |.tmp| is ignored and never executed.
Ordinary FEM/ANS imports and every export retain their companion-file guards.
Import does not launch a native solver or take another step. Capturing the
imported result as a time state permits continuation on its retained mesh.
The native format does not recover the original TikZ source, motion template or
initial source-scaling baseline. Use |.lft| for that full restart contract.
Time problems cannot be exported as an ordinary |.fem| without losing their
history; such an export is rejected. Native control-file generation, transient
AGE exchange and axisymmetric transient ANS conversion are separate limitations.

\subsection{Discrete equations and verification}
For a planar scalar potential, $\mathbf B=(\partial_y A,-\partial_x A)$.
The weak magnetic equation uses the existing nonlinear constitutive law and
adds the consistent P1 conductivity mass matrix:
\[
 \int_\Omega \mathbf H(\mathbf B_{n+1})\cdot\operatorname{curl}v\,dS
 +\int_\Omega\sigma\frac{A_{n+1}-A_n}{\Delta t}v\,dS
 =\int_\Omega(J_s+\sigma e_z)v\,dS+\hbox{boundary terms}.
\]
An unconstrained solid has $e_z=0$. For an isolated planar region, the uniform
$e_z$ enforces zero integrated current; eliminating it gives a symmetric Schur
complement. Each such region is limited to 2048 participating nodes. Stranded
source regions have no resolved conductivity mass. Axisymmetric assembly uses
the existing radial quadrature, azimuthal potential and full-revolution measures;
isolated planar-section constraints do not apply to naturally closed azimuthal rings.

In a moving planar bulk, the difference follows the same material nodes. It
therefore discretises the material derivative and includes the longitudinal
motional $\mathbf v\times\mathbf B$ contribution. No interpolation of history
through a newly generated conducting mesh is performed. A band collision or
quality failure rejects the step before a state is registered.

For a fixed, source-free, linear model with homogeneous boundaries, the discrete
identity is
\[
 W_{n+1}-W_n+\Delta t\,P_{\rm Joule}
 +\tfrac12(A_{n+1}-A_n)^\mathsf{T}K(A_{n+1}-A_n)=0,
\]
with $K$ in physical energy units. The last term is backward-Euler numerical
dissipation. Reversible nonlinear laws give the corresponding convex-energy
inequality. Driven or moving models also exchange source or mechanical work;
do not apply the source-free identity to them unchanged.

Validation includes an analytical diffusion mode, this discrete identity,
axisymmetric/nonlinear dissipation, isolated-current conservation, exact restart
continuation and rotational damping. A time-step refinement gives the expected
first-order trend. Independent FEMM 5 solves on the same mesh agree within
$2.4\times10^{-12}$ relative nodal error for the tested linear steps. Same-mesh
agreement checks implementation; spatial and temporal convergence still need
separate studies for the intended geometry and observables.
